\documentclass[10pt,a4paper]{article}
\usepackage[utf8]{inputenc}
\usepackage[magyar]{babel}
\usepackage[T1]{fontenc}
\usepackage{amsmath}
\usepackage{amsfonts}
\usepackage{amssymb}
\usepackage{makeidx}
\usepackage{graphicx}
\usepackage[pdfusetitle]{hyperref}
\usepackage{titling}
\usepackage{scrextend}

\hypersetup{
    colorlinks,
    citecolor=black,
    filecolor=black,
    linkcolor=black,
    urlcolor=black
}

\date{\today}
\author{Petrányi Bálint}
\title{%
	\textbf{Analízis 2.} \\
	\textbf{Programtervező informatikus A. szakirány} \\
	\medskip
	Definíciók és tételek\\
	\large 2022-2023. tanév 1. félév
}
\setcounter{secnumdepth}{0}
\newcommand{\R}{\mathbb{R}}
\newcommand{\D}{\mathcal{D}}
\newcommand{\fR}{f\in\R\rightarrow\R}
\newcommand{\pf}{+\infty}
\newcommand{\nf}{-\infty}

\begin{document}

\maketitle
\textbf{Remélem nem írtam el semmit illetve hogy jól értelmeztem a kérdéseket ha mégis írjatok hogy mit hibáztam és kijavítom.}

\tableofcontents
\newpage



\section{2. Gyakorlat elméleti kérdései}
\paragraph{1.}
Mi a belső pont definíciója? \\
\\
Legyen $\emptyset \neq A \subset \mathbb{R}$. Azt mondjuk, hogy az $a \subset 	\mathbb{R}$ pont az $A$ halmaz egyes belső pontja, ha $\exists \delta > 0$ olyan, hogy $k_{\delta} \subset A$. \\ \\
Emlékeztető: $k_{\delta}(a)=(a-\delta , a+\delta)$ \\ \\
jelölés: az $A$ belső pontjainak halmazát az $A$ belsejének hívjuk és int $A$-val jelöljük. 

\paragraph{2.}
Mikor mondjuk azt, hogy egy $f \in \mathbb{R} \rightarrow \mathbb{R}$ függvény differenciálható valamely $a \in$ int $\mathcal{D}_f$ \\
\\ 
Akkor ha
\[
\exists \lim_{h \rightarrow 0}{ \frac{f(a+h)-f(a)}{h}} \in \mathbb{R}
\]
Jelölése: $f \in \mathcal{D} \{a\}$

\paragraph{3.}
Mi a kapcsolat a pontbeli differenciálhatóság és a folytonosság között?\\
\\
\[
f \in \mathbb{R} \rightarrow \mathbb{R}, \quad f \in D\{a\} \quad \Longrightarrow \quad  f \in C\{a\}
\]
Szóban: Ha egy valós-valós függvény differenciálható egy pontban, akkor folytonos abban a pontban.
\paragraph{4.}
Milyen ekvivalens átfogalmazást ismer a pontbeli deriválhatóságra lineáris közelítéssel? \\
\\
Legyen $f \in \mathbb{R} \rightarrow \mathbb{R}, a \in \text{int } \mathcal{D}_f$. Ekkor
\[
\begin{split}
f \in & \; D\{a\} \\
& \Updownarrow \\
\exists A \in \mathbb{R}, \; \text{és} \; \epsilon : \mathcal{D}_f \rightarrow \mathbb{R}, \lim_a{\epsilon} = 0 \quad \text{úgy, hogy } & \; f(x)- f(a) = A(x-a)+\epsilon (x)(x-a)
\end{split}
\]
\paragraph{5.}
Mi az érintő definíciója? \\ \\
Legyen $f \in D\{a\}$. Ekkor az 
\[
e_af:\mathbb{R} \rightarrow \mathbb{R}, \quad	 e_af(x) =f'(x)(x-a)+f(a)
\]
egyenesest az $f$ függvény $a$ pontbeli érintőjének hívjuk .
\paragraph{6.}
Milyen tételt ismer két függvény szorzatának valamely pontbeli differenciálhatóságáról és a deriváltjáról? \\ \\
\[
f,g \in D\{a\} \quad \Longrightarrow \quad f \cdot g \in D\{a\}, \; \text{és} \; (f \cdot g)'(a) =f'(a)g(a)+f(a)g'(a)
\]
\paragraph{7.}
Milyen tételt ismer két függvény hányadosának valamely pontbeli differenciálhatóságáról és a deriváltjáról? \\ \\
\[
f,g \in D\{a\}  \; \text{és} \; g(a) \neq 0 \quad \Longrightarrow \quad \frac{f}{g} \in D\{a\}, \; \text{és} \; \Big( \frac{f}{g} \Big) '(a) = \frac{f'(a)g(a)-f(a)g'(a)}{g^2(a)}  
\]
\paragraph{8.}
Milyen tételt ismer két függvény kompozíciójának valamely pontbeli differenciálhatóságáról és a deriváltjáról? \\ \\
Legyen $f,g \in \mathbb{R} \rightarrow \mathbb{R}$, Tegyük fel, hogy az $a$ int $\mathcal{D}_g$ pontban $g \in D\{a\}$, továbbá $f\in D\{g(a)\}$. Ekkor:
\[
f \circ g \in D\{a\}, \; \text{és} \; (f \circ g)' (a) = f'(g(a)) \cdot g'(a)
\]
\paragraph{9.}
Mi az exp, sin, cos függvények derivált függvénye? \\

\begin{itemize}
\item $(e^x)' = e^x$
\item $\sin ' = \cos$
\item $\cos ' = -\sin$ 
\end{itemize}
\newpage

\section{3. Gyakorlat elméleti kérdései}
\paragraph{1.}
Írja fel az $\exp, \; \text{ln}x,\; \sin,\; \cos,\; \text{tg},\; a^x (a>0, x\in\mathbb{R})$ függvények derivált függvényét.  \\
\begin{itemize}
\item $(e^x)' = e^x$
\item $\text{ln }x'=\frac{1}{x}$
\item $\sin ' = \cos$
\item $\cos ' = -\sin$ 
\item $\text{tg }x' = \frac{1}{\cos^2 x }$
\item $ a^{x'} = a \cdot \text{ln }a $
\end{itemize}
\paragraph{2.}
Adjon példát olyan függvényre, ami az $a \in \mathbb{R}$ pontban folytonos, de nem differenciálható! \\
\\
Az abszolút érték függvény folytonos a 0 pontban, de
nem differenciálható 0-ban.
\paragraph{3.}
Mi a jobb oldali derivált definíciója?
\\ \\
Legyen $f\in\mathbb{R}\rightarrow\mathbb{R}, a \in \mathcal{D}_f$ és tegyük fel hogy az $f$ függvény jobbról deriválható (differenciálható), ha 
\[
\exists \quad \text{és véges a} \quad \lim\limits_{x\rightarrow a+0} \frac{f(x)-f(a)}{x-a} \text{határérték}
\] 
Ezt a határértéket nevezzük az $f$ függvény $a$ pontbeli jobb oldal deriváltjának és $f'_+(a)$-val jelöljük 
\paragraph{4.}
Mi az érintő definíciója? \\ \\
Legyen $f \in D\{a\}$. Ekkor az 
\[
e_af:\mathbb{R} \rightarrow \mathbb{R}, \quad	 e_af(x) =f'(x)(x-a)+f(a)
\]
egyenesest az $f$ függvény $a$ pontbeli érintőjének hívjuk .
\paragraph{5.}
Írja le az inverz függvény differenciálszámításáról szóló tételt!
\\ \\
Legyen $I\in\mathbb{R}$ nyílt intervallum és $f: I \rightarrow \mathbb{R}$.
Tegyük fel, hogy
\begin{enumerate}
\item $f$ szigorúan monoton és folytonos $I$-n,
\item $f$ differenciálható az $a \in I$ és $f'(a)\neq 0$.
\end{enumerate} 
Ekkor az $f^{-1}$ inverz függvény deriválható a $b:=f(a)$ pontban, és
\[
\Big( f^{-1} \Big) '(b) = \frac{1}{f'(a)} = \frac{1}{f'(f^{-1}(b))}
\]  
\paragraph{6.}
Milyen tételt hatványsor összegfüggvényének differenciálhatóságáról és a deriváltjáról? \\ \\
Tegyük fel, hogy a $\sum (\alpha_n(x-a)^n) \; (x\in\mathbb{R})$ hatványsor $R$ konvergenciasugara pozitív. Legyen
\[
f(x):= \sum\limits^{+\infty}_{n=0} \alpha_n(x-a)^n \quad (x\in K_R(a))
\] 
a hatványsor összegfügevénye. \\
Ekkor minden $x\in K_R(a)$ pontban $f\in D\{x\}$ és 
\[
f'(x) = \sum\limits^{+\infty}_{n=1} n\alpha_n(x-a)^{n-1} \quad (\forall x \in K_R(a))
\] 
\newpage
\section{4. Gyakorlat elméleti kérdései}
\paragraph{1.}
Milyen szükséges és elégséges feltételt ismer differenciálható függvény monoton növekedésével kapcsolatban?\footnote{(Ha nem jó valaki javítson ki mert nem tudom ere mi lenne a pontos válasz)\label{refnote}} \\ \\
Legyen $I \subset \R $ nyílt intervallum. Tegyük fel hogy $f : I \rightarrow \R, \; f\in D(I)$\\ Ekkor:
\[
f\nearrow \quad \Longleftrightarrow \quad f'\geq 0.
\] 
\textbf{Megjegyzés:} az $f'\geq 0$ feltétel azt jelenti, hogy minden $x\in I$ pontban $f'(x)\geq 0$, aminek geometriai interpretációja az, hogy az érintő meredeksége minden pontban
nem negatív. 
\paragraph{2.}
Milyen elégséges feltételt ismer differenciálható függvény szigorú monoton növekedésével kapcsolatban?\footref{refnote}\\\\
Legyen $I \subset \R $ nyílt intervallum. Tegyük fel hogy $f : I \rightarrow \R, \; f\in D(I)$\\ Ekkor:
\[
f'>0 \quad \Rightarrow \quad f\uparrow
\]
\textbf{Megjegyzés:} Az állítás fordítottja nem igaz. Az $f(x)=x^3 \; (x\in\R)$ függvény szigorúan monoton növekedő, de $f'(0)=0$
\paragraph{3.}
Mit ért azon hogy az $f\in\R\rightarrow\R$ függvény valamely helyen lokális minimuma van? \\\\
Az  $\fR$ függvénynek az $a\in \text{int } \D_f$ pontban lokális minimuma van ha 
\[
\exists K(a), \quad \text{hogy} \quad \forall x \in K(a)\cap \D_f \quad \text{esetén} \quad f(x)\geq f(a)
\]
Ekkor az $a\in\D_f$ pontot az $f$ lokális minimumhelyének nevezzük, $f(a)$ pedig az $f$ lokális minimuma.
\paragraph{4.}
Hogyan szól a lokális szélsőértékre vonatkozó elsőrendű szükséges feltétel?\\\\
Tegyük fel, hogy az $f:\R\rightarrow\R$ függvénynek az $a\in\D_f$ pontban lokális szélsőértéke van és $f\in D\{a\}$ \\
Ekkor:
\[
f'(a)=0
\]
\paragraph{5.}
Adjon  példát olyan $\fR$ függvényre, amelyre valamely $a\in\R$ esetén $f\in D\{a\} , f'(a)=0$ teljesül, de az $f$ függvénynek az $a$ pontjában nincs lokális szélsőértéke \\\\

$f(x)=x^3 \quad (x\in\R)$  esetén $f'(x)=3x^2$ derivált csak a $0$ pontban $0$ ami viszont egyértelműen nem lokális szélsőértékhely $x<0$ esetén $x^3 < 0$ és $x>0$ esetén $x^3>0$

\paragraph{6.}
Mit ért azon, hogy egy függvény valamely helyen jelet vált? \\\\
Azt mondjuk hogy $\fR$ függvény az $a \in \text{int }\D_f$ pontban (-,+) előjelet vált ha $f(a)=0$, és van olyan $\delta >0$, hogy
\[
f(x)<0 \: (a-\delta <x < a), \quad f(x)>0 \: (a<x<a+\delta)
\]
A (+,-) jelváltás értelem szerűen  definiálható
\paragraph{7.}
Hogyan szól a lokális maximumra vonatkozó elsőrendű elégséges feltétel? \\\\
Legyen $\fR$, és a $a\in \text{int }\D_f$.\\
Tegyük fel hogy $\exists \delta >0$, amelyre $f\in D((a-\delta ,a+\delta )) $ és $f'$ előjelet vált $a$-ban \\
Ekkor $f$-nek szigorú lokális szélsőértéke van az $a$-ban \\
($+,-$) jelváltás esetén maximum. \\
\textbf{Megjegyzés:} ($-,+$) jelváltás esetén minimum,
\paragraph{8.}
Írja le a lokális minimumra vonatkozó másodrendű elégséges feltételt. \\\\
Legyen $f \in \R \rightarrow \R$, $a\in \text{int } \D_f$\\
Tegyük fel hogy:
\begin{enumerate}
\item $f\in D^2\{a\}$
\item $f'(a)=0 \quad \text{és} \quad f''(a)\neq 0$  
\end{enumerate}
Ekkor az $a$ pont a szigorú lokális szélsőértékhelye az $f$. függvénynek $f''(a)>0$ esetén minimum \\
\textbf{Megjegyzés:} $f''(a)<0 $ esetén maximum
\paragraph{9.}
Mondja ki a Lagrange-féle középértéktételt! \\
Legyen $a,b\in\R$ és $a<b$. Tegyük fel hogy 
\begin{enumerate}
\item $f\in C [a,b]$
\item $f\in D(a,b)$
\end{enumerate}
Ekkor:
\[
\exists \varepsilon \in (a,b) \quad \text{hogy} \quad f'(\varepsilon) = \frac{f(b)-f(a)}{b-a}
\]
\paragraph{10.}
Mondja ki a Cauchy-féle középértéktételt! \\
Legyen $a,b\in\R$ és $a<b$. Tegyük fel hogy 
\begin{enumerate}
\item $f,g\in C[a,b]$
\item $f,g\in D(a,b)$
\item $g'(x)\neq 0 (x\in (a,b))$
\end{enumerate}
Ekkor:
\[
\exists \varepsilon \in (a,b), \quad \text{hogy} \quad \frac{f'(\varepsilon)}{g'(\varepsilon)} = \frac{f(b)-f(a)}{g(b)-g(a)}
\]
\newpage
\section{5. Gyakorlat elméleti kérdései}
\paragraph{1.}
Definiálja az inflexiós pont fogalmát. \\
Legye $I$ nyílt intervallumon, $\fR, I \subset \D_f$ \\
Azt mondjuk hogy az $a\in I$ pont az $f$ függvénynek inflexiós pontja ha:
\begin{gather*}
\exists \delta > 0 \; k_\delta (a) \subset I \text{ olyan hogy} \\
f \text{ konvex az } (a-\delta , a] \text{ intervallumon és }\\ \text{konkáv az } [a,a+\delta)\text{-n intervallumon vagy fordítva}  
\end{gather*} 
\paragraph{2.}
Mondja ki az inflexiós pont létezésére vonatkozó másodrendű szükséges feltételt. \\
Legye $\fR, \; a \in \text{int } \D_f$ \\
Ha $f\in C^2\{a\}$ és $f$-nek az a pontjában inflexiója van akkor $f''(a)=0$
\paragraph{3.}
Mikor mondjuk, hogy egy függvénynek aszimptotája van a $+\infty$-ben? \\
Legyen $a\in R $ és $f : (a,\pf) \rightarrow \R.$ Azt mondjuk hogy $f$-nek van aszimptotája $(\pf)$-ben, ha:
\[
\exists l(x)=Ax+B \quad (x\in\R)
\]
elsőfokú függvény, amelyre:
\[
\lim\limits_{x\rightarrow\pf}(f(x)-l(x)) = 0
\] 
Ekkor az $l(x) \: (x\in\R)$ egyenes az $f$ aszimptotája $(\pf)$-ben
\paragraph{4.}
Hogyan szól a $(\pf)$-beli aszimptota létezésére vonatkozó tétel?\\
Az $f:(a,\pf)\rightarrow\R$ függvénynek akkor és csak akkor van aszimptotája $(\pf)$-ben, ha léteznek és végesek az alábbi határértékek:
\[
\lim\limits_{x\rightarrow\pf} \frac{f(x)}{x} =: A \in \R \quad \lim\limits_{x\rightarrow\pf} (f(x)-Ax) =: B \in \R
\]
Ekkor az
\[
l(x)=Ax+B \quad (x\in\R)
\]
Egyenes az $f$ függvény aszimptotája $(\pf)$-ben
\paragraph{5.}
Írja le a jobboldali határérték $\frac{0}{0}$ esetére vonatkozó L'Hospital-szabályt. \\
Legyen $\nf \leq a < b < \pf $ és $f,g\in D(a,b).$ Tegyük fel hogy:
\begin{itemize}
\item $\exists \lim\limits_{a+0}f=\lim\limits_{a+0}g=0 $
\item $g(x)\neq$ és $g'(x)\neq 0 \quad \forall x \in (a,b)$
\item $\exists \lim\limits_{a+0} \frac{f'}{g'} \in \overline{\R}$
\end{itemize}
Ekkor
\[
\exists 	\lim\limits_{a+0} \frac{f}{g} \quad \text{ és } \quad \lim\limits_{a+0}\frac{f}{g} = \lim\limits_{a+0}\frac{f'}{g'} \in \overline{\R}
\]
\paragraph{6.}
Írja le a baloldali határérték $\frac{\pf}{\pf}$ esetére vonatkozó L'Hospital-szabályt. \\
\end{document} 